\documentclass[10pt,a4paper]{amsart}
\usepackage{amsmath,amssymb,amsthm,graphicx}
\usepackage{amsmath}
\usepackage{amssymb}
\usepackage{hyperref}
\usepackage[utf8]{inputenc}
\usepackage{csquotes}
\usepackage[foot]{amsaddr}


\usepackage{times}
\usepackage{amsmath,amsthm,amssymb}
\usepackage{array}
\usepackage{bbm}
\usepackage{xcolor}
\usepackage{verbatim}
\usepackage{url}
%\usepackage{authblk}
\pagestyle{plain}


\def\ve{\varepsilon}
\def\vf{\varphi}
\def\vn{\varnothing}

\def\AA{\mathcal{A}}
\def\AAA{\mathbb{A}}
\def\BB{\mathcal{B}}
\def\CC{\mathcal{C}}
\def\CCC{\mathbb{C}}
\def\DDD{\mathbb{D}}
\def\DD{\mathcal{D}}
\def\FF{\mathcal{F}}
\def\HH{\mathcal{H}}
\def\II{\mathcal{I}}
\def\III{\mathbb{I}}
\def\MM{\mathcal{M}}
\def\NN{\mathcal{N}}
\def\NNN{\mathbb{N}}
\def\PPP{\mathbb{P}}
\def\QQQ{\mathbb{Q}}
\def\RR{\mathcal{R}}
\def\RRR{\mathbb{R}}
\def\SS{\mathcal{S}}
\def\SSS{\mathbb{S}}
\def\UU{\mathcal{U}}
\def\VV{\mathcal{V}}
\def\WW{\mathcal{W}}
\def\ZZ{\mathcal{Z}}
\def\ZZZ{\mathbb{Z}}
\def\PP{\mathcal{P}}
\def\TTT{\mathbb{T}}
\def\EEE{\mathbb{E}}


%%%%  jb abbreviations
\def\bK{{\bf K}}
\def\bL{{\bf L}}
\def\bU{{\bf U}}
\def\aut{\mathop{aut}}
\def\betabar{{\bf\beta}}
\def\Uscr{\mathcal U}
\def\myreduct{\upharpoonright}
\def\abar{{\bf a}}
\def\bbar{\bf b}
\def\cbar{\bf c}
\def\lg{\mathop{\rm lg}}
\def\nmdeg{\mathop{\rm nmdeg}}
\def\icl{\mathop{\rm icl}}
% \sidebar{..} puts .. in comment form
\newcommand{\sidebar}[1]{\vskip10pt\noindent
 \hskip.70truein\vrule width2.0pt\hskip.5em
 \vbox{\hsize= 4truein\noindent\footnotesize\relax #1 }\vskip10pt\noindent}



 %%%%

\def\QQQa{\mathbb{Q}\langle \Bar{\alpha},\alpha_0 \rangle}

\def\cc{\mathfrak{c}}

\def\a{\alpha}
\def\b{\beta}
\def\g{\gamma}
\def\d{\delta}
\def\e{\varepsilon}
\def\z{\zeta}
\def\l{\lambda}
\def\k{\kappa}
\def\o{\theta}
\def\r{\rho}
\def\s{\sigma}
\def\t{\tau}
\def\w{\omega}
\def\D{\Delta}

\def\nmdeg{\operatorname{nmdeg}}
\newcommand{\Fmb}[1]{F_{Mb}\left( #1  \right)}
\newcommand{\indf}[1]{F_{ind}\left( #1  \right)}
\newcommand{\fcb}[1]{F_{cb}\left( #1  \right)}
\newcommand{\ficb}[1]{F_{Icb}\left( #1  \right)}





\DeclareMathOperator{\alg}{alg}
\DeclareMathOperator{\RU}{RU}
\DeclareMathOperator{\RM}{RM}
\DeclareMathOperator{\RD}{RD}
\DeclareMathOperator{\tp}{tp}
\DeclareMathOperator{\acl}{acl}
\DeclareMathOperator{\dcl}{dcl}
\DeclareMathOperator{\Cb}{Cb}
\DeclareMathOperator{\dtr}{\delta-trdeg}


\def\id{\operatorname{id}}
\def\Frac{\operatorname{Frac}}
\def\Const{\operatorname{Const}}
\def\spec{\operatorname{Spec}}
\def\span{\operatorname{span}}
\def\exc{\operatorname{Exc}}
\def\Div{\operatorname{Div}}
\def\cl{\operatorname{cl}}
\def\mer{\operatorname{mer}}
\def\trdeg{\operatorname{trdeg}}
\def\ord{\operatorname{ord}}




\newtheorem{theorem}{Theorem}[section]
\newtheorem{thm}[theorem]{Theorem}
\newtheorem{question}[theorem]{Question}
\newtheorem{lemma}[theorem]{Lemma}
\newtheorem{cor}[theorem]{Corollary}
\newtheorem{prop}[theorem]{Proposition}
\newtheorem{fact}[theorem]{Fact}
\newtheorem{claim}[theorem]{Claim}
\newtheorem{subclaim}[theorem]{Subclaim}
\newtheorem{conj}[theorem]{Conjecture}
\newtheorem{cond}[theorem]{Condition}
\newtheorem{observation}[theorem]{Observation}
\newtheorem{notation}[theorem]{Notation}

\theoremstyle{definition}
\newtheorem{definition}[theorem]{Definition}
\newtheorem{example}[theorem]{Example}
\newtheorem{xca}[theorem]{Exercise}
\newtheorem{ques}[theorem]{Question}
\newtheorem{note}[theorem]{Note}

\theoremstyle{remark}
\newtheorem{remark}[theorem]{Remark}
\newcommand{\m}{\mathbb }
\newcommand{\mc}{\mathcal }
\newcommand{\mf}{\mathfrak }
\newcommand{\mb}{F_{Mb} }
\newcommand{\G}{\mathrm{G} }
\newcommand{\U}{\mathrm{U} }



\def \ta {\tau_{\mathcal{D}/\Delta}}
\def \D {\Delta}
\def \DD {\mathcal D}

\newcommand{\gen}[1]{\left\langle#1\right\rangle}
%\renewcommand{\restriction}{\mathord{\upharpoonright}}

\newcommand{\pd}[2]{\frac{\partial #1}{\partial #2}}
\newcommand{\td}{\text{tr.deg.}}
\newcommand{\pp}{\partial }
\newcommand{\pdtwo}[2]{\frac{\partial^2 #1}{\partial #2^2}}
\newcommand{\od}[2]{\frac{d #1}{d #2}}
\def\Ind{\setbox0=\hbox{$x$}\kern\wd0\hbox to 0pt{\hss$\mid$\hss} \lower.9\ht0\hbox to 0pt{\hss$\smile$\hss}\kern\wd0}
\def\Notind{\setbox0=\hbox{$x$}\kern\wd0\hbox to 0pt{\mathchardef \nn=12854\hss$\nn$\kern1.4\wd0\hss}\hbox to 0pt{\hss$\mid$\hss}\lower.9\ht0 \hbox to 0pt{\hss$\smile$\hss}\kern\wd0}
\def\ind{\mathop{\mathpalette\Ind{}}}
\def\nind{\mathop{\mathpalette\Notind{}}}
\numberwithin{equation}{section}
\newcommand{\is}{^{p^ {-\infty}}}
\newcommand{\QQ}{\mathbb Q}
\newcommand{\fh}{\mathfrak h}
\newcommand{\SL}{\operatorname{SL}}


\title{Simple Homogeneous Structures and Indiscernible Sequence Invariants}
\author[1]{John T. Baldwin}
\author[2]{James Freitag}
\author[3]{Scott Mutchnik}
%\affil[1]{Department of Mathematics, Statistics, and Computer Science, University of Illinois at Chicago}
%%%\affil[2]{Department of Biology, University Y}
%%\author{ John T. Baldwin}
%\author{James Freitag} 
\address{Department of Mathematics, Statistics, and Computer Science, University of Illinois at Chicago }
\email{jbaldwin@uic.edu}
\email{jfreitag@uic.edu}

%\author{Scott Mutchnik}
%\address{Department of Mathematics, Statistics, and Computer Science, University of Illinois at Chicago}
\email{mutchnik@uic.edu}

\thanks{JB was partially supported by Simons travel grant G3535. JF was partially supported by NSF CAREER award number 1945251. SM, who is additionally affiliated with the Institut de mathématiques de Jussieu – Paris Rive Gauche, was supported by the NSF under Grant No. DMS-2303034, and by the European Union’s Horizon 2020 research and innovation programme under the Marie Sklodowska-Curie grant agreement No 101034255.}

\date{\today}
\usepackage{cleveref}
\crefname{originalremark}{original remark}{original remarks}
\Crefname{originalremark}{Original remark}{Original remarks}
\makeatletter
\newlabel{findbigger}{{2.1.5}{source}{}{sourcepointer-1}{}}
\newlabel{findbigger@cref}{{[originalremark][0][]2.1.5}{[1][0][]source}{}{}{}}
\makeatother
\title{In differentially closed fields the non-minimality degree agrees with the fibre dimension}
\author{John Baldwin, James Freitag, Scott Mutchnik}
\date{arxiv:2405.08211v3; CC BY 4.0}
\begin{document}
\maketitle
\noindent\textit{Editorial note.} All delivered mathematical text is the author's own, copied from the pinned publisher source sources/logic-fl03-simple-homogeneous-structures-and/pinned-original-source.tex; no statement or proof step was rewritten, reordered or removed.\par\medskip
\noindent\textit{Editorial note.} Preamble provenance: the retained preamble is the source's own layout adapter, if any, plus the author's macro block at lines 3-205 of sources/logic-fl03-simple-homogeneous-structures-and/pinned-original-source.tex (everything between the publisher class metadata and \textbackslash{}begin\{document\}, comments excluded). Nothing was added, no macro was invented, and the publisher document class is deliberately not reproduced; the excerpt serves only as preamble, never as a span coordinate.\par\medskip
\noindent\textit{Editorial note.} Metadata provenance: source\_id and source\_version were taken from sources/logic-fl03-simple-homogeneous-structures-and/license-observation.json; the title from SOURCE.json; the author list from the source metadata record; the source URL from the first non-empty of source\_url, source\_urls[0] and the evidence's fixed\_version\_url/license\_url. The licence evidence file used by this card is sources/logic-fl03-simple-homogeneous-structures-and/license-observation.json, recorded in license.evidence\_path so that the binding target is explicit.\par\medskip
\noindent\textit{Editorial note.} External references are navigation pointers only; every deliberately kept step of the author's argument stays inside the counted spans. The displayed numbers of those pointers are the author's own numbering under the author's \textbackslash{}newtheorem declarations (section.index as declared), computed mechanically by the lane's scanner.\par\medskip
\noindent\textit{Editorial note.} Bibliography provenance (Chief-approved route, 2026-10-12): this source directory stages no bibliography at all - neither a thebibliography block nor a .bib database - while the counted proof cites four keys. The citations are therefore resolved with the BibTeX database staged in the ACL lane's own copy of the same paper, focused-seeds/ACL/sources/simple-homogeneous-2405.08211v3/refs.bib, whose licence evidence records source\_id arxiv:2405.08211v3 and source\_version 2405.08211v3 - the same arXiv version this card's pinned original was read from. All four keys (moosa2014some, knop1983mehrfach, freitag2023bounding, freitag2023degree) were verified present in that database, and no entry was edited in any way; the backend is bibtex, not a hand-made thebibliography.\par\medskip
\noindent\textit{Editorial note.} The counted proof contains one reference to another labelled item, the remark labelled findbigger, recorded as a navigation stub of kind remark.\par\medskip
\noindent\textit{Allowed standard prerequisites (editorial).}\begin{itemize}
\item The theory DCF\_0 of differentially closed fields of characteristic zero, and the two invariants compared: the non-minimality degree and the fibre dimension F\_ind. Set up in the paragraph preceding the counted statement (source lines 440-441). Fixes the theory and the two quantities whose equality is asserted.
\item The classification of types with no proper fibrations into locally modular and non-locally-modular cases. Cited inside the counted proof (Moosa, Proposition 2.3) and used at lines 446-451. Supplies the dichotomy that splits the proof into its two cases.
\item Knop's classification of multiply transitive algebraic group actions. Cited inside the counted proof at line 455. Identifies the binding group action with affine transformations or the action of PGL on projective space, which pins the fibre dimension to 2.
\end{itemize}\medskip
\ifdefined\hypertarget\hypertarget{sourcepointer-1}{}\fi
\noindent\textit{Source reference.} Original remark 2.1.5: sources/logic-fl03-simple-homogeneous-structures-and/pinned-original-source.tex, source lines 451--451. Navigation pointer for \textbackslash{}ref\{findbigger\} used inside the counted span (line 451); the argument that consumes it is retained in the span.\par\medskip
\begin{prop} \label{dcfsame}
    For a type $p$ in the theory $DCF_0$, $\nmdeg(p) = F_{ind} (p).$
\end{prop}

\begin{proof}
Let $p \in S(A)$. If $p_n \rightarrow p_{n-1} \rightarrow \ldots \rightarrow p_1$ is an analysis of the type $p$ over $A$, $F_{ind} (p)$ is bounded by $F_{ind} (p_1)$ and also $F_{ind}$ of the types given by the generic fibers of the maps in the analysis. Thus, it suffices to prove the result in the case that $p$ has \emph{no proper fibrations} \cite[Definition 2.1]{moosa2014some}. It follows by \cite[Proposition 2.3]{moosa2014some} that $p$ is either 
\begin{enumerate}
    \item interalgebraic over $A$ with a finite collection of independent realizations of a locally modular minimal type, or 
    \item interalgebraic with a type $q$ which is internal to a nonlocally modular type. 
\end{enumerate}
    In the first case, both $F_{ind} (p)= \nmdeg{(p)} =1$. In the second case, if the binding group of $p$ relative to the constants is not generically $2$-transitive \cite[for the defintion of generic transitivity]{freitag2023bounding}, then $F_{ind}(p) = 1$, and $\nmdeg(p) \leq F_{ind} (p) $ in general (see Remark \ref{findbigger}), so the quantities are equal. In the second case, since $F_{ind} (-) $ is invariant up to interalgebraicity, we can assume $p$ is internal to the nonlocally modular type - in this case, the constants. 

    Suppose now that $F_{ind} (p) >1$. Then it must be that $p$ and also $p \otimes p$ are weakly orthogonal to the constants over $A$ (if not, it is easy to see that $F_{ind} (p) =1$ by taking a Morley sequence whose image in $\mc C^n$ coming from nonorthogonality lies on an algebraic curve). It follows that both $p$ and $p \otimes p$ are isolated types. We wish now to reduce to the case that the action of the binding group $G$ on $p$ is $2$-transitive. This is equivalent to any two realizations of $p$ being independent. As the binding group is generically $2$-transitive, with $p \otimes p$ being isolated, there must be an $A$-definable equivalence relation $\sim$ on $p$ such that $p/ \sim$ has a $2$-transitive binding group. 

    But now the binding group is $2$-transitive and definably isomorphic to an algebraic group action acting on an algebraic variety over a field of characteristic zero, and so the classification theorem of Knopp \cite{knop1983mehrfach} tells us that the binding group action on $p$ is isomorphic to either a group of affine transformations of affine space or the action of $PGL_{n+1}$ on $\m P^n.$ Here it must be that $n>1.$ Both of these group actions preserve collinearity of points. Taking the Morley sequence of points on a line, we can see that $F_{ind} (p) =2.$ $\nmdeg(p)=2$ in this case as well since the binding group acts $2$-transitively. 
\end{proof}

\bibliographystyle{plain}
\bibliography{refs}
\end{document}
